\kafkasection{A bright spot of plasma in the jet of a quasar is observed to have moved between 2 observations separated by 10 years. The quasar is at a distance of 120 Mpc from Earth. The displacement between the 2 observations was measured to be 0.04 arcseconds. }


\kafkasubsection{how fast is its apparent motion on the sky (in cm/s) and expressed as v/c?}
At a distance of $d = \qty{120}{Mpc}$, the physical distance across the sky would be:
\begin{equation}
    r = \theta d = (\qty{0.04}{''})(\qty{206265}{^{-1}''/rad})(\qty{120}{Mpc})
\end{equation}
\begin{equation}
    = \qty{23.3}{pc}
\end{equation}
Expressed in cgs we have:
\begin{equation}
    r = \qty{7.18e19}{cm}
\end{equation}
Over the course of 10 years, we get:
\begin{equation}
    \dot{r} = \frac{\qty{7.18e19}{cm}}{10\cdot\qty{3.154e7}{s}} = \qty{2.28e11}{cm/s}
\end{equation}
Uh oh that's $\qty{7.59}{c}$
\kafkasubsection{What angle is the plasma moving with respect to our line-of-sight to the quasar (in degrees) if gamma=1000?
Hint: Use Figure 1.5 and related equations of special relativity to solve part ii.
}
The angle is given by:
\begin{equation}
    \theta = \frac{1}{\gamma} = \qty{0.001}{rad}
\end{equation}
